% !TEX root = ./main.tex

\chapter{相关技术与理论基础}

AIGC安全分析涉及生成模型、伪造检测、内容审核和平台工程实现。本章围绕深度伪造与多模态内容生成方法，虚假内容检测，不安全内容审核与判定，以及前后端分离、代理调度和模型服务化等工程基础展开，为后续系统设计提供依据。

\section{深度伪造与多模态内容生成方法}

\subsection{图像生成技术演进}

图像生成技术主要经历了从生成对抗网络到扩散模型的演进。GAN通过生成器与判别器之间的对抗训练学习数据分布，在人脸生成、图像翻译、风格迁移和属性编辑等任务中较早取得突破\cite{ref1,ref2,ref3}。其生成速度较快，但训练稳定性不足，复杂场景和高分辨率细节表达存在局限。

扩散模型通过正向加噪和逆向去噪完成生成建模，在训练稳定性和图像质量上更具优势\cite{ref4,ref5}。以Latent Diffusion Model和Stable Diffusion为代表的模型在潜空间中扩散和去噪，降低了计算成本，并提升了文本对齐、高分辨率表达和可控生成能力\cite{ref6,ref7,ref19}。从安全视角看，生成质量提升使早期依赖GAN指纹、频域伪影和颜色异常的检测方法不再稳定，检测器需要更强的跨模型泛化能力。

扩散模型的工程化应用通常由文本编码器、去噪网络、采样调度器和图像解码器共同完成。采样步数、提示词权重、随机种子和后处理策略都会影响图像质量，也会改变纹理、边缘和噪声表现。因此，AIGC检测不能只依赖某一种固定伪影，而需要同时考虑模型来源、生成参数和后处理链路对检测特征的影响。

\subsection{视频生成技术路线}

视频生成不仅要求单帧视觉质量，还要保持时间维度上的运动连续性和语义一致性。常见路线包括文本到视频、图像到视频和人物驱动动画。文本到视频模型根据文本条件生成连续画面；图像到视频模型以首帧或参考图为基础生成动态片段；人脸或人体驱动模型则利用关键点、运动场或驱动视频迁移表情、口型和姿态。

以本文原型系统接入的ModelScope Text-to-Video类模型为例，其通常采用扩散机制和时空建模结构生成短视频\cite{ref20}。这类模型能够生成较自然的片段，但在复杂动作、长时一致性、文字渲染和物理规律表达上仍可能出现局部闪烁、结构漂移和运动不连续等问题，这些现象也为视频检测提供了线索。

与图像生成相比，视频生成的安全分析需要额外关注时间维度。单帧画面可能具备较高真实感，但相邻帧之间仍可能出现身份漂移、边界跳变、纹理闪烁和运动方向不一致等问题。对于人脸或数字人视频，还需要关注口型、眼部动作、头部姿态和语音语义是否协调。因此，视频检测不仅需要视觉分类，还需要对整段内容的动态语义进行综合判断。

\subsection{面向人脸的深度伪造技术}

Deepfake是AIGC生成技术在身份迁移和人物驱动场景中的典型形态，广义上指利用深度学习模型对人脸身份、表情、口型、语音或姿态进行替换、重建和驱动的技术集合\cite{ref13,ref14,ref21,ref22}。身份替换关注将源身份迁移到目标场景；表情迁移和口型同步关注动作与语音的一致性；属性编辑则通过文本或标签修改年龄、发型、妆容和背景等特征。它本质上仍属于AI生成或编辑内容，只是风险集中体现在人物身份与社会关系上，因此不宜与AIGC生成技术完全割裂。

人脸伪造的风险后果更直接，容易触发身份冒用、名誉损害、舆论误导和诈骗等问题。因此，人物图像和人物视频既是虚假检测的重要对象，也是系统设计中必须覆盖的典型输入类型。

从检测角度看，人脸伪造既有局部线索，也有跨帧线索。局部线索包括面部边缘融合不自然、肤色与光照不一致、眼镜和头发等细节区域变形；跨帧线索则包括眨眼频率异常、口型与语音不匹配、面部身份在运动中轻微漂移等。由于这些线索会随生成算法和后处理方式变化，本文平台没有将Deepfake检测限定为单一人脸算法，而是把它纳入更一般的图像与视频虚假内容检测流程。

\subsection{面向通用场景的多模态伪造技术}

AIGC伪造风险已从人脸换脸扩展到通用场景构造。文生图模型能够快速生成新闻场景、灾害画面和拟真肖像，文本到视频和图像到视频模型还能生成带有动作和镜头变化的动态内容。与人脸伪造相比，通用场景伪造更难依赖固定局部区域检测，往往涉及物理规律错误、对象关系异常、时序动作不自然和语义逻辑不一致等问题。因此，检测方法需要更强的开放域泛化能力。

通用场景伪造还具有明显的跨模态特征。生成者不仅可以通过文本提示构造图像和视频，还可以结合参考图、局部重绘、背景替换、风格迁移、字幕和配音生成更完整的叙事内容。经过二次剪辑、平台压缩和社交传播后，伪造内容往往不再保留清晰的模型痕迹，单纯依靠人脸区域或单帧局部异常难以完成判断。因此，平台需要同时关注图像、视频、文字和场景语义之间的一致性。

生成内容通常在低层统计、局部语义、时序动态和安全对齐方面留下可分析特征。低层统计特征包括噪声分布、纹理连续性、频谱结构和压缩残留差异；局部语义特征包括边缘不稳定、细粒度结构错误和区域不连贯；视频时序特征包括运动轨迹异常、帧间差异不自然和长时结构漂移；安全对齐问题则表现为模型在诱导性提示下可能输出违规内容。

这些特征并不总是稳定出现。高质量生成模型会削弱明显伪影，社交平台压缩、裁剪和滤镜也可能掩盖原始痕迹；真实图像在低光照、强压缩或快速运动下也可能产生类似异常。因此，平台不能把单一线索等同于最终结论，而应通过置信度、风险标签和原因说明表达不确定性。

因此，本文同时关注两类任务：一类识别内容是否由AI生成或被深度伪造，另一类识别内容是否具有不安全属性。前者强调真实性和来源分析，后者强调风险等级和治理要求，二者在平台中分别建模、统一展示。

\section{虚假内容检测方法}

\subsection{基于低层统计与频域特征的检测方法}

较早的伪造检测方法依赖图像处理和媒体取证思想，通过颜色分布、噪声模式、压缩残留、插值痕迹和频域特征识别异常\cite{ref9,ref15,ref22}。这类方法可解释性较强，适合辅助定位异常区域，但对后处理、压缩、裁剪和分辨率变化较敏感。随着扩散模型生成质量提高，早期明显伪影被削弱，因此低层统计方法更适合作为辅助分析，而难以单独支撑开放场景检测。

\subsection{基于监督分类的检测方法}

监督分类方法将伪造检测建模为二分类或多分类任务，利用卷积神经网络、Transformer等视觉编码器从真实与伪造样本中学习判别边界\cite{ref9,ref11}。在人脸伪造场景中，FaceForensics++等数据集推动了面部裁剪、局部注意力、压缩鲁棒性和多任务学习等研究\cite{ref13}。这类方法在训练集和测试集分布相近时效果较好，但面对未见过的新模型、压缩链路或跨模态场景时容易退化。

监督分类方法的工程优势在于推理流程明确，模型输出通常可以直接转换为概率、标签和阈值判定，便于服务化部署和批量评测。但其效果高度依赖训练样本的覆盖范围。如果训练集中主要包含某一类GAN图像，模型可能学习到该生成器特有的纹理痕迹，而不是更通用的“AI生成”特征。因此，在本文平台中，本地模型更适合作为可复验基线和实验入口，而不应被理解为覆盖所有开放场景的唯一判断来源。

\subsection{基于通用视觉表征与重建对比学习的检测方法}

为缓解跨模型泛化不足，研究者开始利用大规模预训练视觉模型或视觉语言模型提取通用特征，再通过近邻搜索、线性分类头或轻量探测器判断真实性\cite{ref8}。UniversalFakeDetect就是该思路的代表，其借助预训练特征空间中的判别能力提高对未知生成模型的适应性\cite{ref10}。

重建与对比学习方法则通过构造高质量重建样本和困难伪造样本，引导检测器关注更稳定的生成痕迹。以DRCT为代表的研究表明，若模型能区分与真实图像高度相似的困难伪造样本，在未知扩散模型上也更可能保持性能\cite{ref12}。这类方法更符合生成模型快速迭代背景下的检测需求。

通用表征方法与本文系统设计关系较密切。一方面，它们不要求平台重新训练大型视觉主干，适合在毕业设计阶段形成可运行的本地检测服务；另一方面，预训练特征本身仍需要阈值选择、结果解释和异常处理。也就是说，算法模型只能给出判别分数，平台层还必须负责把分数转换为用户可理解的结论、置信度和说明文本。

\subsection{面向视频的时序检测方法}

视频检测比图像检测多出时间维度，因此更关注帧间一致性、动作规律和动态差分。传统方法可逐帧检测后再聚合，也可使用三维卷积、循环网络或时空Transformer建模短片段。近年来，训练自由或弱训练方法开始关注真实视频与AI生成视频在运动变化规律上的差异。以D3为代表的方法利用时间差分和二阶动态特征检测AI生成视频\cite{ref16}，说明单帧逼真并不代表整体时序自然。

在进入综合讨论前，本文先将上述主要检测路线进行归纳，便于后文说明平台为何采用本地模型与外部多模态能力并行的方案。相关比较如表~\ref{tab:fake-detection-methods} 所示。

\begin{table}[H]
\centering
\caption{虚假内容检测方法比较}
\label{tab:fake-detection-methods}
\renewcommand{\arraystretch}{1.35}
\zihao{5}\songti
\begin{tabular}{C{3.4cm}L{3.2cm}L{3.2cm}L{3.4cm}}
\toprule
\textbf{方法类别} &
\multicolumn{1}{C{3.2cm}}{\textbf{主要依据}} &
\multicolumn{1}{C{3.2cm}}{\textbf{优势}} &
\multicolumn{1}{C{3.4cm}}{\textbf{局限与适用场景}} \\
\midrule
低层统计与频域方法 & 噪声、纹理、频谱、压缩痕迹 & 可解释性较强，适合媒体取证辅助分析 & 对后处理敏感，面对高质量扩散模型时泛化不足 \\
监督分类方法 & 标注数据中学习到的判别边界 & 封闭数据集内精度较高，部署链路成熟 & 依赖训练分布，跨生成模型和跨压缩链路容易退化 \\
通用视觉表征方法 & 大规模预训练模型特征空间 & 泛化能力较好，适合作为开放域图像检测基线 & 对预训练特征质量敏感，通常仍需阈值或轻量分类头 \\
视频时序检测方法 & 帧间一致性、运动规律、动态差分 & 能利用单帧图像中不明显的时序异常 & 计算和数据处理成本较高，本地部署复杂度更大 \\
多模态大模型判别 & 视觉输入与自然语言提示联合推理 & 接入灵活，适合开放场景解释与快速原型 & 成本、稳定性和可解释依据仍需平台层面约束 \\
\bottomrule
\end{tabular}
\end{table}

\subsection{不同检测方法的比较与适用场景}

由表~\ref{tab:fake-detection-methods} 可见，不同检测路线各有适用场景：低层统计方法解释性强但泛化不足；监督分类方法封闭集效果好但依赖训练分布；通用视觉表征方法更适合开放图像检测；视频时序方法能利用动态异常但部署复杂；多模态大模型接入灵活、解释能力较强，但成本、稳定性和依据透明度仍需约束。

对于本文平台，单一方法难以同时覆盖图像、视频、真实性和安全性任务。因此，本文采用“本地图像检测模型 + 外部多模态理解能力”的混合方案：本地模型提供可控、可复验的研究基线，外部接口补充视频理解、开放场景判别和安全审核能力。

\section{不安全内容检测方法}

\subsection{审核标签体系与判定目标}

内容安全审核关注内容是否违反平台规则、法律规范或伦理要求，而不是判断其是否由AI生成。在AIGC场景下，生成模型可能输出暴力、色情、仇恨、毒品、赌博或其他违规内容，也可能产生边界模糊的擦边、暗示和符号化风险。因此，审核任务首先需要定义风险标签体系，再围绕每类标签寻找可观察证据，最后给出通过、复审或阻断等处置建议。

不安全内容不是由单个视觉目标直接决定的，而是由目标、动作、场景和上下文共同判定。例如，检测到“刀具”并不必然意味着暴力风险，若同时出现攻击动作、受伤对象或血迹，风险等级才会明显升高；检测到人物身体区域也不必然属于色情风险，还需要结合裸露程度、姿态、遮挡关系和画面语境；敏感符号、违法活动和赌博等风险则常依赖文字、标志、道具和场景组合判断。由此可见，分类和目标检测只是提取证据的基础步骤，真正的审核结论还需要风险规则或多模态语义推理进行归纳。

表~\ref{tab:unsafe-evidence} 归纳了常见风险类别、可观测证据与判定逻辑。后续系统实现中，平台并不直接暴露底层检测器的全部细节，而是将外部审核接口输出统一映射为风险标签、风险等级、处置建议和原因说明。

\begin{table}[H]
\centering
\caption{不安全内容审核的证据与判定逻辑}
\label{tab:unsafe-evidence}
\renewcommand{\arraystretch}{1.35}
\zihao{5}\songti
\begin{tabular}{C{2.5cm}L{4.8cm}L{5.9cm}}
\toprule
\textbf{风险类别} &
\multicolumn{1}{C{4.8cm}}{\textbf{可观测证据}} &
\multicolumn{1}{C{5.9cm}}{\textbf{判定逻辑}} \\
\midrule
暴力血腥 & 武器、攻击动作、血迹、伤口、倒地对象、打斗场景 & 单一武器展示通常进入复审；若与攻击动作、伤害结果或血腥画面同时出现，则提高为高风险。 \\
色情低俗 & 大面积裸露、敏感部位、性暗示姿态、成人场景道具 & 结合人体区域检测、姿态和语境判断；边界模糊内容多进入复审，明确裸露或性行为倾向进入阻断。 \\
仇恨与敏感符号 & 极端标志、歧视性文字、侮辱性手势、群体攻击语义 & 需要图像符号识别、OCR和语义分类共同判断，避免仅因普通图案或无关文字造成误判。 \\
违法活动 & 毒品、赌博器具、非法交易场景、危险行为 & 依据道具识别、场景识别和文字线索综合判定，若出现明确交易、使用或诱导语义则提高风险等级。 \\
正常内容 & 未出现上述高风险目标、动作、文字或语境 & 若分类置信度低且无明显风险证据，则给出低风险或通过建议；若证据不足但存在疑似线索，则转入复审。 \\
\bottomrule
\end{tabular}
\end{table}

\subsection{基于分类、检测与OCR的图像审核}

传统图像审核通常由多标签分类、目标检测、OCR和规则聚合组成。多标签分类模型从整张图像判断是否接近暴力、色情、赌博等风险类别；目标检测模型定位武器、血迹、人体区域、毒品、赌博器具等局部证据；OCR和符号识别用于提取画面中的文字、标志和敏感符号。三类结果分别提供全局语义、局部区域和文本证据，避免只依赖单一路径。

在判定阶段，系统会把不同证据转换为同一风险类别下的分数。例如，对于某一风险类别$c$，可将图像分类置信度、该类别相关目标的最高检测置信度以及OCR命中得分进行融合：
\[
S_c=\max(P_c^{cls}, \max_i P_{c,i}^{det}, P_c^{ocr})
\]
其中，$P_c^{cls}$ 表示分类模型给出的类别置信度，$P_{c,i}^{det}$ 表示第$i$个相关目标的检测置信度，$P_c^{ocr}$ 表示文字或符号线索命中的风险得分。随后再结合风险规则进行修正：若高风险目标与高风险动作同时出现，则上调等级；若只有单个低置信度线索，则进入复审而不是直接阻断。这样，目标检测和分类不只是“说有风险”，而是为审核结论提供可解释的视觉证据。

\subsection{视频审核中的时序聚合}

视频审核比图像审核多出时间维度，不能简单地把某一帧结果等同于整段视频结论。常见做法是按固定fps或镜头变化抽取关键帧，对每帧执行图像审核，再结合动作识别、帧间变化和持续时间进行聚合。若某类风险只在个别帧中低置信度出现，通常应降低结论强度；若风险目标、动作或文字在多个连续片段中稳定出现，或构成完整的攻击、裸露、交易等过程，则应提高风险等级。

因此，视频不安全内容判定通常需要同时回答三个问题：第一，风险证据在什么时间段出现；第二，风险证据是否持续或重复出现；第三，前后帧是否形成明确动作和语义关系。本文平台的视频审核接口保留fps参数，正是为了控制采样密度，使外部视频理解能力能够在成本和时序覆盖之间取得平衡。

\subsection{基于多模态大模型的结构化安全审核}

多模态大模型能够同时接收图像、视频和文本提示，并结合自然语言规则输出风险结论、标签和解释。与传统分类器相比，其优势在于可以把物体、动作、文字和场景关系放在同一语义空间中分析，适合开放类别和复合风险判断。工程上，审核提示词需要明确风险范围和输出格式，例如要求模型围绕暴力、色情、仇恨、敏感符号、违法活动等类别判断是否安全，并返回 \texttt{safe}、\texttt{suggestion}、\texttt{labels} 和 \texttt{reason} 字段。

多模态大模型的结论仍需要平台层约束。首先，平台应要求其给出具体原因，原因中应包含可观察证据，如“画面中出现武器并伴随攻击动作”或“画面文字包含赌博诱导信息”，而不是只有抽象的“不安全”。其次，后端需要把 \texttt{pass}、\texttt{review} 和 \texttt{block} 等建议映射为低、中、高风险等级，并处理返回格式不稳定、字段缺失和解析失败等异常情况。最后，审核结果应与虚假内容检测分开展示，因为“AI生成”与“违规风险”是两个不同判断维度。

\subsection{生成安全评测与防御研究的启示}

近年来，研究者还关注生成模型本身的安全评测和前置防御。例如，视频生成安全评测基准从多类风险维度评估文本到视频模型安全性，说明视频生成不仅要关注画质和时序一致性，还要关注输出是否触发具体风险类别\cite{ref17}。相关有害性识别数据集也表明，生成安全正在从“输出检测”扩展到“提示词输入、模型生成、结果审核和平台处置”的全流程治理\cite{ref18}。本文平台主要聚焦生成后检测与分析，因此在设计上需要同时覆盖真假鉴别、内容合规、风险证据说明和处置建议表达。

\section{系统开发关键技术}

\subsection{前端交互与可视化技术}

面向AIGC安全分析的平台不仅要完成模型调用，还要提供清晰稳定的交互界面。React与TypeScript适合构建模块化前端，Vite提供轻量开发和构建能力，Ant Design提供表单、上传、标签、进度和反馈等常用组件。前端展示层的核心职责是将复杂模型调用转化为统一操作流程，包括文件上传、链接输入、任务选择、进度反馈、结果可视化和历史记录查看。

\subsection{后端代理与服务编排技术}

本文涉及本地模型、外部审核接口和多种输入输出协议，后端需要承担统一代理和服务编排功能。Node.js与Express具有实现成本低、异步处理方便和前端生态衔接自然等特点。后端代理层负责鉴权封装、参数校验、格式转换、任务路由、异步轮询、错误处理和结果标准化，是实现多源能力统一组织的关键。

\subsection{模型服务化部署技术}

本地检测和生成功能需要以服务化方式接入平台。FastAPI等Python Web框架便于将PyTorch模型封装为HTTP接口，完成模型加载、输入预处理、推理、后处理和状态监测。模型服务化能够将算法脚本转化为可复用组件，使前端、后端和模型之间形成清晰边界，也为模型替换、远程部署和多模型扩展提供基础。

对本文系统而言，服务化部署使本地模型、外部接口和前端页面以统一协议协同工作。模型侧关注预处理和推理结果，后端侧负责鉴权、路由、容错和字段映射，前端侧处理交互和可视化展示。这种分层方式降低单个模块变化对整体系统的影响，也为后续接入更多检测模型和审核接口提供工程基础。

\section{本章小结}

本章梳理了深度伪造与多模态内容生成方法、虚假内容检测方法、不安全内容审核与判定方法以及平台实现相关工程技术。综合来看，生成模型演进使伪造内容来源更复杂、逼真程度更高，而检测与审核方法各有优势和局限。面向真实应用场景的AIGC安全平台需要采用多源能力协同、统一结果结构和模块化部署相结合的设计思路。
